% 12_getting_started_yc.tex — How to Get Started & Doing Things That Don't Scale by YC Founders

\section{How to Get Started \& Doing Things That Don't Scale by YC Founders (Stanley Tang and Walker Williams)}
\label{sec:getting_started_yc}

Stanley Tang is the founder of Doordash, an online food delivery company that relies on individual drivers fulfilling mobile orders instead of a centralized structure.
He shares insights from the early days of Doordash and three key tips when starting ambitious projects: \textbf{testing hypotheses, launching fast, and doing things that ``do not scale''.}

The Doordash co-founders wanted to build technology for small businesses.
They created a landing page called PaloAltoDelivery.com as an experiment to gauge demand (``testing hypotheses'').
They built the original website in one hour, and despite offering a suboptimal user experience and barely doing any marketing, orders came in quickly (``launching quickly'').
The \textbf{need for their product was validated} by the continued consumer demand even without specialized user design.

In the early days, the founders relied on improvised strategies such as using Apple's Find My Friends app to track their delivery drivers, taking orders through their personal phone number, and even making deliveries themselves (``doing things that don't scale'').
After joining Y Combinator, they were able to incorporate their company and leverage more efficient technologies in order to grow their business, but Tang argues that the most important thing at first is simply \textbf{getting off the ground.}

Walker Williams is the founder of Teespring, an e-commerce company that allows entrepreneurs to launch apparel brands.
He expands on Tang's idea of doing things that do not scale --- things that are fundamentally unsustainable --- in the context of three strategies for early startups: \textbf{finding first users, turning those users into champions, and finding product-market fit.}

\textbf{Unfortunately, there is no magical strategy for user acquisition.}
It is difficult for founders to sell a new product because they do not deeply understand their users' pain points yet.
As a result, Williams advocates that it is a founder's responsibility to \textbf{do whatever it takes to bring in their first customers,} which usually includes a lot of personal effort such as emails or phone calls.
The benefit of finding initial users is growth, not an immediate return on investment.

\textbf{A champion is a user who advocates for the product.}
Williams says that the best way to gain champions is to give users good experiences that they will \textit{remember}, and one good (but eventually unsustainable) way to do so is to constantly talk to them.
For instance, Williams himself handles dozens of customer service emails every day, proactively reaches out to current (and ``churned'') customers, and reads through social media comments and other online communities.
Since problems are inevitable, founders must also \textbf{do whatever it takes to make those issues right.}

Similarly to Tang, Williams cautions that the initial product and features will not be the product and features that the company eventually grows with, by virtue of talking to more customers and learning their true pain points.
\textbf{It is crucial to optimize for speed over scalability or clean code.}
